\begin{workedexample}[Worked Example: When the sample dominates noise]
A receive coil's noise is the sum of coil (conductor) and sample (body)
losses. The ratio of unloaded to loaded quality factor reveals which dominates:
\[ \frac{Q_{\text{unloaded}}}{Q_{\text{loaded}}} = 1 + \frac{R_{\text{sample}}}{R_{\text{coil}}}. \]
A ratio of 5 means the sample resistance is $4\times$ the coil resistance --- the
sample dominates, so polishing the conductor buys little. At low field or for
small coils the coil can dominate, and conductor quality matters more.
\end{workedexample}

\begin{keytakeaways}
\begin{itemize}[leftmargin=1.4em,itemsep=2pt]
\item Total noise $\propto R_{\text{coil}} + R_{\text{sample}}$; SNR is best when the sample dominates.
\item $Q_{\text{unloaded}}/Q_{\text{loaded}} = 1 + R_{\text{sample}}/R_{\text{coil}}$ diagnoses the regime.
\item Sample noise rises with field; small or low-field coils are more coil-noise-limited.
\end{itemize}
\end{keytakeaways}

\begin{exercises}
\item If $Q_{\text{unloaded}}/Q_{\text{loaded}} = 3$, what is $R_{\text{sample}}/R_{\text{coil}}$?
\item In the sample-dominated regime, does better conductor plating help much?
\item Are small surface coils more coil- or sample-noise-limited?
\end{exercises}

\furtherreading{Mispelter~\cite{mispelter}; Roemer \emph{et al.}~\cite{roemer}.}
